% Source: Wald GR, physical PDF pp. 66--67
%         (printed pp. 53--54).
% Coverage: Chapter 3 problems 1--8.
% Figures/tables/footnotes: no figures, tables, or footnotes.
% Status: source-reviewed and compile-clean.
% Uncertainties: none.

\chapterproblems

\begin{enumerate}[label=\arabic*.,leftmargin=*]
\item Let property (5) (the \lq\lq torsion free\rq\rq\ condition) be dropped from the definition of
derivative operator $\nabla_a$ in section 3.1.
  \begin{enumerate}[label=\alph*),leftmargin=2em]
  \item Show that there exists a tensor $T^c{}_{ab}$ (called the \emph{torsion tensor}) such that for all
  smooth functions, $f$, we have $\nabla_a\nabla_bf-\nabla_b\nabla_af=-T^c{}_{ab}\nabla_cf$.
  (Hint: Repeat the derivation of Eq.~\eqref{eq:3.1.8}, letting $\widetilde\nabla_a$ be a torsion-free
  derivative operator.)
  \item Show that for any smooth vector fields $X^a$, $Y^a$ we have
  \[
    T^c{}_{ab}X^aY^b=X^a\nabla_aY^c-Y^a\nabla_aX^c-[X,Y]^c.
  \]
  \item Given a metric, $g_{ab}$, show that there exists a unique derivative operator $\nabla_a$ with
  torsion $T^c{}_{ab}$ such that $\nabla_cg_{ab}=0$. Derive the analog of Eq.~\eqref{eq:3.1.29}, expressing
  this derivative operator in terms of an ordinary derivative $\partial_a$ and $T^c{}_{ab}$.
  \end{enumerate}

\item Let $M$ be a manifold with metric $g_{ab}$ and associated derivative operator $\nabla_a$. A
solution of the equation $\nabla_a\nabla^a\alpha=0$ is called a \emph{harmonic function}. In the case where
$M$ is a two-dimensional manifold, let $\alpha$ be harmonic and let $\epsilon_{ab}$ be an antisymmetric
tensor field satisfying $\epsilon_{ab}\epsilon^{ab}=2(-1)^s$, where $s$ is the number of minuses occurring
in the signature of the metric. Consider the equation $\nabla_a\beta=\epsilon_{ab}\nabla^b\alpha$.
  \begin{enumerate}[label=\alph*),leftmargin=2em]
  \item Show that the integrability conditions (see problem 5 of chapter 2 or appendix
  B) for this equation are satisfied, and thus, locally, there exists a solution, $\beta$. Show
  that $\beta$ also is harmonic, $\nabla_a\nabla^a\beta=0$. ($\beta$ is called the harmonic function
  \emph{conjugate} to $\alpha$.)
  \item By choosing $\alpha$ and $\beta$ as coordinates, show that the metric takes the form
  \[
    ds^2=\pm\Omega^2(\alpha,\beta)\bigl[d\alpha^2+(-1)^sd\beta^2\bigr].
  \]
  \end{enumerate}

\item
  \begin{enumerate}[label=\alph*),leftmargin=2em]
  \item Show that $R_{abcd}=R_{cdab}$.
  \item In $n$ dimensions, the Riemann tensor has $n^4$ components. However, on account
  of the symmetries~\eqref{eq:3.2.13}, \eqref{eq:3.2.14}, and~\eqref{eq:3.2.15}, not all of these components are
  independent. Show that the number of independent components is $n^2(n^2-1)/12$.
  \end{enumerate}

\item
  \begin{enumerate}[label=\alph*),leftmargin=2em]
  \item Show that in two dimensions, the Riemann tensor takes the form
  \[
    R_{abcd}=R\,g_{a[c}g_{d]b}.
  \]
  (Hint: Use the result of problem 3(b) to show that $g_{a[c}g_{d]b}$ spans the
  vector space of tensors having the symmetries of the Riemann tensor.)
  \item By similar arguments, show that in three dimensions the Weyl tensor vanishes
  identically; i.e., for $n=3$, Eq.~\eqref{eq:3.2.28} holds with $C_{abcd}=0$.
  \end{enumerate}

\item
  \begin{enumerate}[label=\alph*),leftmargin=2em]
  \item Show that any curve whose tangent satisfies Eq.~\eqref{eq:3.3.2} can be reparameterized
  so that Eq.~\eqref{eq:3.3.1} is satisfied.
  \item Let $t$ be an affine parameter of a geodesic $\gamma$. Show that all other affine
  parameters of $\gamma$ take the form $at+b$, where $a$ and $b$ are constants.
  \end{enumerate}

\item The metric of Euclidean $\mathbb R^3$ in spherical coordinates is
\[
  ds^2=dr^2+r^2(d\theta^2+\sin^2\theta\,d\phi^2)
\]
(see problem 8 of chapter 2).
  \begin{enumerate}[label=\alph*),leftmargin=2em]
  \item Calculate the Christoffel components $\Gamma^\sigma{}_{\mu\nu}$ in this coordinate system.
  \item Write down the components of the geodesic equation in this coordinate system
  and verify that the solutions correspond to straight lines in Cartesian coordinates.
  \end{enumerate}

\item As shown in problem 2, an arbitrary Lorentz metric on a two-dimensional manifold locally always can be put in the form $ds^2=\Omega^2(x,t)[-dt^2+dx^2]$. Calculate
the Riemann curvature tensor of this metric (a) by the coordinate basis methods of
section 3.4a and (b) by the tetrad methods of section 3.4b.

\item Using the antisymmetry of $\omega_{\lambda\mu\nu}$ in $\mu$ and $\nu$, Eq.~\eqref{eq:3.4.15}, show that
\[
  \omega_{\lambda\mu\nu}
  =3\omega_{[\lambda\mu\nu]}-2\omega_{[\mu\nu]\lambda}.
\]
Use this formula together with Eq.~\eqref{eq:3.4.23} to solve for $\omega_{\lambda\mu\nu}$ in terms of
commutators (or antisymmetrized derivatives) of the orthonormal basis vectors.
\end{enumerate}
